% !TEX program = lualatex
\documentclass[11pt,a4paper]{ltjsarticle}
\usepackage{amsmath,amssymb,amsthm,mathtools,luatexja-fontspec}
\usepackage[margin=25mm]{geometry}
\usepackage[colorlinks=true,urlcolor=blue,linkcolor=blue]{hyperref}
\setmainjfont{HaranoAjiMincho-Regular}
\newtheorem{theorem}{定理}[section]
\newtheorem{proposition}[theorem]{命題}
\newtheorem{lemma}[theorem]{補題}
\newtheorem{corollary}[theorem]{系}
\newtheorem{remark}[theorem]{注意}
\newcommand{\CC}{\mathbb C}
\newcommand{\PP}{\mathbb P}
\newcommand{\Sym}{\operatorname{Sym}}
\newcommand{\MF}{\operatorname{MF}}
\newcommand{\Perf}{\operatorname{Perf}}
\renewcommand{\proofname}{証明}
\title{Sklyanin混合族の有限模型と形式比較}
\date{2026年11月13日}
\begin{document}
\maketitle
\begin{abstract}
混合Sklyanin族の点イデアル、評価射、積のhomotopyを有限行列で記述する。
曲率付き有限arity gaugeを実際の次数付き行列因子化へ移す仕組みを明示し、
非冪零な対象微分への無限和を用いない。
係数を微分しない有限縮約と相対windowを組み合わせ、積とSerre traceから混合一次単元を全次数で決定する。
全二十一混合方向を六級数へ縮め、標識付きHesse平面incidenceとの形式同定を得る。
有限非零量子係数における圏比較の代数化は主張しない。
\end{abstract}
\section{点イデアルと混合欠損}
$R=\CC[[h]]$、$S=R\langle x,y,z\rangle/([x,y]-hz^2,[y,z]-hx^2,[z,x]-hy^2)$とする。
$D=x^3+y^3+z^3$、$T=\sum_{\pi\in S_3}\pi(xyz)$、$G=D-h^2T/3$、$\Omega=G/3$と置く。
中心性と形式PBW性は現行の中心三次対論文\cite{Pair}の証明による。
\begin{proposition}\label{mix:point}
$l=x+y,Q=x^2-xy+y^2$とし、
\begin{align*}
 a&=z+h(y-x),&b&=-l-h^2z,\\
 q_1&=Q+hz^2-h^2xz,&q_2&=z^2+hxz-2hyz+h^2Q,\\
 r_1&=z^2-hxz-h^2y^2,&r_2&=-Q-hz^2
\end{align*}
と置く。次の行列は$G$および$\Omega$の行列因子化を与える：
\[
 \phi=\begin{pmatrix}a&q_1\\b&q_2\end{pmatrix},\quad
 \psi=\begin{pmatrix}r_1&r_2\\l&z\end{pmatrix},\quad
 P=\phi\operatorname{diag}(1,1/3),\quad
 Q^{\rm MF}=\operatorname{diag}(1,3)\psi/3.
\]
$P_0=S(-1)^2,P_1=S(-2)\oplus S(-3)$として$P:P_1\to P_0$、
$Q^{\rm MF}:P_0(-3)\to P_1$である。
$C=S/(G)$上の余核は点$o=[1:-1:0]$のイデアル
$I_o=\ker(C\to S/(lS+zS))$である。
\end{proposition}
\begin{proof}
$la+zb=hr_{xy}+r_{yz}-r_{zx}=0$である。
残る八成分は自由語として展開し、二次関係の左右倍の線形包で零になることを確認する。
現行検証器は$\mathbb Q(h)$上で重み四以下の全関係倍を構成し、各成分を厳密に簡約する。
$h$を打ち切らない恒等式である。
複体$0\to S(-2)\xrightarrow{(a,b)^t}S(-1)^2\xrightarrow{(l,z)}S$は、
$h=0$で通常のKoszul複体なので、各内部次数の単元pivotを持ち上げると完全性と余核の階数一の自由性を得る。
$\phi\psi=\psi\phi=G$から$G$はその余核を零化する。
$C$上の$(l,z)$の核元を$S^2$へ持ち上げ、その像$Gc$を$(q_1,q_2)^tc$で引けば、
$S$上の核になり$(a,b)^td$と書ける。従って$\operatorname{coker}\phi=I_o$である。
最後の基底変更は表示通りで、周期性の係数$1/3$を保持する。
\end{proof}

楕円曲線$E=V(D+h^2xyz)$、$o=[1:-1:0]$、$p=[1:-1:-h]$、$q=-p=[1:-1:h]$とする。
右twisted homogeneous coordinate ringの規約では、window $i=1$で$I_o$は$\mathcal O_E(-o)$を表す。
実際その次数$n\ge1$切断は$H^0(E,\mathcal O_E(-o)\otimes L_n)$であり、次数$3n-1>0$なので高次切断が消える。
Picard tensor $-\otimes\mathcal O_E(p-o)$での像は$\mathcal O_E(-q)$である。
$\tau=3(x+y)-h^2z$について$\operatorname{div}\tau=3o$、
$\operatorname{div}(x+y)=o+p+q$なので、$3(x+y)/\tau$がこの線束同型を与える。
可換側は$k=z-hx,Q_c=Q+h^2xz,R_c=z^2+hxz$と置き
\[
 \phi_c=\begin{pmatrix}k&Q_c\\-l&R_c\end{pmatrix},\qquad
 \psi_c=\begin{pmatrix}R_c&-Q_c\\l&k\end{pmatrix}
\]
である。$lQ_c+kR_c=D+h^2xyz$が行列因子化を与える。
$\Omega$と同じ規格化の$P_c,Q_c^{\rm MF}$には同じ対角基底変更と$1/3$を用いる。

\begin{proposition}[混合欠損と八行模型]\label{mix:defect}
関係$[y,z]=h(x^2+m)$、他の二関係はそのままとし、$m$を重み二の中央多項式変数とする。
\begin{align*}
 R_0&=\begin{pmatrix}x&h(x-2y)/3\\(h^2x+2hz)/3&x+2h^2z/3\end{pmatrix},\\
 R_1&=\begin{pmatrix}x+2h^2z/3&-J\\-h&x\end{pmatrix},\\
 9J&=h^2x^2+2hxz-2h^2yx+4h^2y^2-2hyz-2hzy-2h^3z^2
\end{align*}
とすると
\[
 PQ^{\rm MF}=(\Omega+mx)I-mR_0,\quad Q^{\rm MF}P=(\Omega+mx)I-mR_1,
 \quad R_0P=PR_1,\quad R_1Q^{\rm MF}=Q^{\rm MF}R_0.
\]
$m=\alpha uv$とし、$C_j=u^2I+\alpha vR_j$、
\[
 P_u=\begin{pmatrix}P&uI\\-C_1&Q^{\rm MF}\end{pmatrix},\quad
 Q_u=\begin{pmatrix}Q^{\rm MF}&-uI\\C_0&P\end{pmatrix},\quad
 \Phi=\begin{pmatrix}P_u&vI\\-v^2I&Q_u\end{pmatrix},\quad
 \Psi=\begin{pmatrix}Q_u&-vI\\v^2I&P_u\end{pmatrix}
\]
とすると$\Phi\Psi=\Psi\Phi=(\Omega+\alpha uvx+u^3+v^3)I_8$である。
\end{proposition}
\begin{proof}
四つの二行行列恒等式は自由語の関係空間で検証する。
後半二式は前半、中心性、および形式PBWでの$m$の正則性からも従う。
$PC_1=C_0P,Q^{\rm MF}C_0=C_1Q^{\rm MF}$なのでblock積の非対角項は零となり、
対角は$PQ^{\rm MF}+uC_0$と$Q^{\rm MF}P+uC_1$である。
$v$についてもう一度同じ二因子積を取れば八行模型を得る。
自由加群は$u$追加後$E_0^u=E_0\oplus E_1(2),E_1^u=E_1\oplus E_0(-1)$、
$v$追加後も同じ規則でshiftする。
特殊ファイバーの$\Phi$は$F_0$の行列因子化として単射であり、各内部次数の単元minorを持ち上げると余核は基底上平坦になる。
\end{proof}

\section{評価方向と積の有限homotopy}
分離した$m=0$へ戻る。上の点イデアルMFを$M$とする。
\begin{proposition}
$M\oplus M(2)$は楕円MF因子を生成する。
次の行列対は$M\to M(2)$の次数$-1$の閉射で、そのHom群の自由階数一の生成元である：
\[
 A=\begin{pmatrix}(hy-z)/3&0\\0&1\end{pmatrix},\qquad
 B=\begin{pmatrix}1&h^2x/3\\h^2&(-hx+2hy+(h^3-1)z)/3\end{pmatrix}.
\]
$\Delta_0=3+h^3$として可換側の生成元は
\[
 A_c=\begin{pmatrix}s&t\\-r&1\end{pmatrix},\quad B_c=\operatorname{adj}A_c,
 \quad s=-z/3-2hx/\Delta_0,\ t=(4h^2x-2h^2y)/(3\Delta_0),\ r=2h^2/\Delta_0.
\]
\end{proposition}
\begin{proof}
楕円windowで$M=\mathcal O_E(-o)$であり、二回のrotationは点対象$\mathcal O_{-p}[1]$を与える。
線束と一点のskyscraperは、点での完全列を繰り返して全ての整数twistを生成し、それらはample twistなので$\Perf(E/R)$を生成する。
行列のclosed条件$PA+BQ^{\rm MF}=Q^{\rm MF}B+AP=0$は二次関係の左右倍として厳密に確認できる。
次数$-2$の境界射の可能形は内部次数から$\left(\begin{smallmatrix}0&0\\0&c\end{smallmatrix}\right)$だけなので、
その境界の$A$の右下成分は零である。
表示射の右下成分は一であり、線束から点対象へのHomが自由階数一なので生成元である。
可換側では$Q_c^{\rm MF}=\operatorname{adj}P_c$なのでclosed条件は$\operatorname{tr}(P_cA_c)=0$へ帰着し、直接展開で分かる。
\end{proof}

\begin{proposition}\label{mix:product}
$A_* =\left(\begin{smallmatrix}-1/3&-2h/3\\0&0\end{smallmatrix}\right)$、
$B_* =\left(\begin{smallmatrix}0&2h/3\\0&(h^3-1)/3\end{smallmatrix}\right)$とすると
\[
 PA_*+B_*Q^{\rm MF}=BA+hR_0,\qquad
 Q^{\rm MF}B_*+A_*P=AB+hR_1.
\]
従って指定周期性を通して評価射を二回合成した類は$[f^2]=-h[r]$となる。
\end{proposition}
\begin{proof}
両blockの自由語展開を二次関係の左右倍で簡約する。
現行検証器は混合欠損・整合式・評価cocycle・この二境界式をすべて$\mathbb Q(h)$で独立に検証する。
$[r]$は混合変分の積と曲率を同時に微分して得るcharacteristic類であり、単に乗法$x$を取った類ではない。
この点は混合欠損を微分した式$\dot d\,d+d\dot d+\dot\mu(d,d)=\dot W I$にも現れる。
\end{proof}

\section{係数複体と対象評価で失われるcoherence}
\subsection{楕円回転核}
$f:E\to\operatorname{Spec}R$を上の楕円族、$i:\Delta\hookrightarrow E\times_RE$を対角とする。
$A_P=-\otimes P$、$\rho_P=A_P\rho_0A_P^{-1}$と置く。
Fourier--Mukaiの規約を$\Phi_K(F)=Rp_{2*}(p_1^*F\otimes^{\mathbf L}K)$とすると
\[
 K_P=(p_1^*P^{-1}\otimes p_2^*P)\otimes I_\Delta\otimes p_1^*L[1]
\]
が回転核である。$I_\Delta$はCartier因子の可逆イデアル層であり、二項分解そのものではない。
\begin{proposition}
自然な擬同型
\[
 \operatorname{RHom}_{E\times_RE}(\mathcal O_\Delta,K_P)
 \simeq Rf_*L\simeq Rx\oplus Ry\oplus Rz
\]
があり、右辺は次数零に集中する。
指定周期性$\rho_P^3\simeq[2]$で負二回転を読むと
$HC_R^\bullet(\Perf(E),\rho_P^{-2})\simeq(Rx\oplus Ry\oplus Rz)[-2]$となる。
全てのArtin基底変換と両立する。
\end{proposition}
\begin{proof}
Cartier双対性$i^!M=Li^*M\otimes T_{E/R}[-1]$を用いる。
対角上の$P^{-1}\otimes P$は標準評価で消え、
$i^*I_\Delta=\Omega^1_{E/R}$なので$i^!K_P=L$である。
これは局所自由なCartier分解でも計算でき、Artin基底の非平坦性は障害にならない。
$0\to\mathcal O_{\PP^2_R}(-2)\to\mathcal O_{\PP^2_R}(1)\to i_{E*}L\to0$から
$Rf_*L=R^3[0]$である。
共役による正一回転の写像はこの基底上恒等だが、指定周期性の写像まで恒等になるとは限らない。
\end{proof}

グラフだけで核を書くと、基底上の直線を落とす危険がある。
$\Gamma_p=\{y=x+p\}$、差写像$d(x,y)=y-x$なら$I_{\Gamma_p}=d^*\mathcal O_E(-p)$であり、
裸のグラフ核による係数複体には
\[
 \Lambda_p=o^*\mathcal O_E(-p)\otimes_R\operatorname{Lie}(E/R)
\]
がtensorされる。局所環上で自明でも、そのframeには$1+hR$の単元差がある。
本族の形式座標$t=z/x$で$p$のイデアルgeneratorは$t+h$であり、
その$o$への制限と$\partial_t$のtensorが特殊ファイバーで一となるframeを与える。
これは可逆層のgeneratorであって数値$h$そのものではない。

正一回転比較を$\eta$、周期性defectを$\kappa$とし$c^3=\kappa^{-1}$とすると、
文字通りの負二回転比較は$c^{-2}\eta^{(-2)}$となる。
指定周期性で$\rho^{-2}=\rho[-2]$と読めば
$\eta^{(-2)}=\kappa^{-1}\eta[-2]$なので係数は
$c^{-2}\kappa^{-1}=c$である。
従って正一回転基底で測る混合係数は$cq$であり、$q=1$を仮定してはならない。

\subsection{一変数A2因子}
定数族$\MF^{\mathbb Z}(R[u],u^3)$には、二頂点一矢$A_2$の固定dg模型を用いる。
体上の同値とrotationの同定はTakahashi\cite{Takahashi}の既知入力であり、
固定した生成子模型を$R$へ基底変換して用いる。
$P_1=(R\xrightarrow1R),P_2=(0\to R)$、$f:P_2\hookrightarrow P_1$とすると、
$\rho=S^{-1}[1]$は
\[
 \rho P_1=P_2[1],\quad\rho P_2=[P_2\xrightarrow fP_1],\quad
 \rho f:[P_2\to P_1]\to P_2[1]
\]
で表される。最後の写像は次数$-1$で恒等、次数零で零である。
path algebraの長さ一の対角分解を用いると
\[
 HC_R^\bullet(\MF^{\mathbb Z}(u^3),\rho)
 \simeq[R\xrightarrow{c\mapsto(c,c)}R^2]\simeq R[0].
\]
二項の次数は$-1,0$であり、$\pi(a,b)=b-a,i(c)=(0,c),H(a,b)=a$は
$dH+Hd=1-i\pi$を満たす整数係数の分裂を与える。
Jacobian側の掛け算$u$の類でこの自由直線を規格化する\cite{BFK}。
この模型は$h$に依存しないので未知の$h$単元を加える必要はない。

ここで$H^0\operatorname{Hom}(P_i,\rho P_i)=0$だが$HC_R^0=R$である。
生成元$(0,1)$の第二成分は矢$f$に対するcoherenceを表す。
従って生成対象ごとの評価が境界であることからHochschild類を零と結論できない。

三因子の分離Thom--Sebastiani模型とcoherentな$(\mathbb Z/3)^2$作用を固定すれば、
$g=\rho_2\rho_3\rho_1^{-2}$の係数複体は、perfect対角と係数bimoduleのKünneth同型から
$(Rx\oplus Ry\oplus Rz)[-2]$となる。
特殊ファイバーでの三つの類は$uvx,uvy,uvz$であり、内部次数シフトと可換である。
各群元の作用行列$M$は$M^3=I,M\equiv I\pmod h$を満たす。
$M^2+M+I\equiv3I$は可逆なので$M=I$であり、不変部分を取っても階数三は変わらない。
これは係数複体の計算であり、源の混合cocycleをこの空間へ移した具体的写像の計算とは区別する。

\section{有限arity gaugeの行列因子化への移送}
可換$\mathbb Q$代数$K$、$S=K[x,y,z,u,v]$、周期変数
$\deg_{\rm coh}\beta=-2,\deg_{\rm int}\beta=3$を取る。
$n$入力cochainに$\beta^j$を掛けたもののshift後次数を$n-1-2j$とする。
正規化とは正arity成分が入力$1$を持つと零になることである。
arity方向には直和を取り、無制限な積完備化は用いない。
通常のMF $P:E_1\to E_0,Q:E_0(-3)\to E_1$は、$S[\beta,\beta^{-1}]$上の奇微分
\[
 d=\begin{pmatrix}0&P\\\beta^{-1}Q&0\end{pmatrix},\qquad d^2=F/\beta
\]
と等しい。
\begin{lemma}[有限移送]\label{mix:finite}
$I^N=0$となる基底イデアルを取り、二つの曲がった結合構造$M,M'$が$I$を法として一致するとする。
正規化された有限arity cochainの有限和$\xi\in I\,CC^0_{\rm fin}(S[\beta,\beta^{-1}])[1]$が
$K[\beta,\beta^{-1}]$線形、内部次数零で、$M'=e^{\operatorname{ad}\xi}M$を満たせば、
有限自由graded MF圏の間に$I$を法として恒等な$A_\infty$同値がある。
これはcone、有限直和、直和因子へ延長し、指定された三次周期性を保つ。
\end{lemma}
\begin{proof}
bar coderivationの指数$\mathcal F=e^\xi$は$I^N=0$により有限和であり、
各入力cochainも有限arityなのでTaylor成分は有限個である。
共役式は$\mathcal F M=M'\mathcal F$、すなわち曲がった関手恒等式そのものである。
各Taylor成分を行列成分へ拡張し、中間添字を和し、shiftにはKoszul符号を用いる。
対象微分の像は$s d'=\sum_{n\ge0}\mathcal F_n(sd,\ldots,sd)$であり、
射の像も両端と間へ$d$を挿入した有限和で定める。
$s$の次数を$-1$とし、$b_0=s(F/\beta),b_2(sa,sb)=(-1)^{|a|}s(ab)$とすれば、
零曲率twist条件は正確に$d^2=F/\beta$である。
関手恒等式のarity零部分が$(d')^2=F'/\beta$を与える。
$d$自身の冪零性は使わず、全ての和は基底冪零性とarityで止まる。
$e^{-\xi}$にも同じ操作を行えばbraceの合成恒等式により互いに逆になる。
正規化は厳密単位性を与え、$\beta$線形性から
$\mathcal F_1(\beta\operatorname{id})=\beta\operatorname{id}$、
その入力を持つ高次成分は零である。
この式はtwist後も成立するので指定周期性を保つ。
有限twisted complexと冪等完備化への延長は同じ有限行列公式で行う。
\end{proof}

\section{有限多重微分cochainの縮約}
$S_0=K[x,y,z]$で係数を微分しないcochainを用いる。
$D^a=\partial^a/a!$の入力型$D^{a_1}\otimes\cdots\otimes D^{a_n}$に対し、
各$a_i\ne0$とし、全微分指数$A=\sum a_i$を固定する。
Hochschild微分は入力の一つを二つの正多重指数に分ける有限写像
\[
 d_H=\sum_i(-1)^{n+i-1}\sum_{b+c=a_i,\ b,c\ne0}
 (a_1,\ldots,a_{i-1},b,c,a_{i+1},\ldots,a_n)
\]
である。
\begin{lemma}\label{mix:cobar}
arity四以上の閉cochainには有限なarity一つ下の原始がある。
閉二入力cochainは有限一入力境界とHKR双ベクトルの和になる。
いずれの構成も係数に線形で、全微分指数と内部次数を保つ。
\end{lemma}
\begin{proof}
固定$A$の入力分割は、正の真部分多重指数のposet
$\mathcal P_A=\{s:0\le s\le A,\ s\ne0,A\}$の正規化単体鎖に等しい。
ある$A_i\ge2$なら$c(s)=s\vee e_i$は依然このposetに属し、
$s\le c(s),e_i\le c(s)$を満たす。
prism homotopy $P$で$\partial P+P\partial=c_*-\operatorname{id}$、
$c(\mathcal P_A)$を$e_i$へconeする$K$で$\partial K+K\partial=\operatorname{id}$となる。
従って$H=Kc_*-P$は縮約であり、その双対にarity符号$(-1)^{p-1}$を掛ければ表示の微分の原始を得る。
三方向でarity四以上なら$|A|\ge4$なので必ずこの場合に入る。
全て有限posetの有限和で、係数多項式を変えない。

二入力では閉性は正の三分割について$c(a+b,c)=c(a,b+c)$である。
$|A|\ge3$の比較グラフは連結なので全係数が一つの値に等しく、一入力$-c_AD^A$の境界である。
$A=2e_i$も同様で、$A=e_i+e_j$だけは二係数の平均を境界で除くと反対称部分が残る。
これはHKR双ベクトルである。
\end{proof}

\begin{proposition}[全次数での純曲率代表]\label{mix:pure}
$K=\mathbb Q[e_1,e_2,e_3,m_1,m_2,m_3]$、$\operatorname{wt}e_i=1,\operatorname{wt}m_i=2$とし、
$F=D/3+\sum e_ix_jx_k+\sum m_ix_i$と置く。
形式PBW積の各$h$係数が有限な$K$線形多重微分作用素であるとする。
実際の中心曲率を持つ源$M_h$は、各$h^N=0$で有限arityの係数線形gaugeにより
$\mu+P_h/\beta$という純曲率へ移せる。
一つの補空間を固定して、
\[
 P_h=F+\sum_{r\ge1}h^rP_r(x,e),\qquad P_r\in\mathbb Q[x_1,x_2,x_3,e_1,e_2,e_3]_3
\]
と選べる。補正の補空間は五十六次元である。
\end{proposition}
\begin{proof}
$D_0=[\mu+F/\beta,-]$に関して閉じる次数一・内部次数零の有限cochainを考える。
最大arityが四以上なら補題\ref{mix:cobar}でその成分を除く。
曲率とのbracketはarityを一つ下げるので、最高成分を取り直す有限降下でarity二以下になる。
残る双ベクトル$\Pi$は$\iota_{dF}\Pi=0$を満たす。
三勾配$F_i=x_i^2+e_jx_k+e_kx_j+m_i$は、$m_i$を消去できるため多項式環で正則列である。
Koszul完全性から$\Pi=\iota_{dF}\Theta$であり、三ベクトルのHKR像を原始として除ける。
残る零arityは$g/\beta$である。
重み三の$g$を一次ベクトル場の像$\sum S_1F_i$で割ると$m_i$を消去でき、
商は六つの重み一変数$x,e$の三次部分である。その次元は$\binom83=56$。
各段階は有限で係数に線形である。

$h$次数帰納で、低次が通常積と純曲率になっていれば次の誤差は$D_0$に関して閉じる。
純曲率同士のbracketは零なので新たな障害はなく、上の原始を$h^r$倍して指数化すればよい。
固定した補空間・縮約を全切断で用いるとArtin切断間に両立する系になる。
PBW積の有限微分性は還元グラフ表示\cite{BW}から得る。
各$h^r$には有限個の$r$頂点グラフだけが現れ、全入力微分次数は高々$2r$である。
中央係数に向かう交換補正は零なので係数微分は現れない。
\end{proof}

\paragraph{現在の適用範囲.}
純曲率の存在と有限移送の補題だけでは、標準的可換対応での座標や周期性単元の値は定まらない。
相対Artin window、楕円側のrotation剛性、characteristic類とSerre traceの比較を別に確認する必要がある。
有限$h$への収束・代数化は全て別問題である。

\section{Artin基底の相対window}
局所Artin環$R$上、$B$は五変数多項式環の次数別有限自由二次変形、
$W\in Z(B)_3$は滑らかな三次式へ還元するとする。
相対対角が長さ五の有限自由Koszul分解を持ち、左右で
$\operatorname{RHom}_B(R,B)=R(5)[-5]$と仮定する。
混合族では相対PBW性と有限自由Frobenius二次双対がこの仮定を与える。
$W$は左右正則で、$A=B/(W)$も次数別有限自由とする。
冪零filtrationと特殊ファイバーのNoether性から$A,B$は左右Noetherである。
\begin{theorem}\label{mix:window}
\[
 D_{\rm rel}(\operatorname{gr}A)
 =\{M\in D^b(\operatorname{gr}A):M|_B\in\Perf(\operatorname{gr}B)\},\qquad
 D_{\rm sg,rel}(A)=D_{\rm rel}/\Perf(\operatorname{gr}A)
\]
と置く。有限自由graded MFの冪等完備化は$D_{\rm sg,rel}(A)$と$R$線形Morita同値である。
実際の斉次Ore局所化による構造代数層を$\mathcal A$とすると、同じwindowと二つのmutationにより
\[
 D_{\rm sg,rel}(A)\simeq
 \mathcal K(\mathcal A)=\langle\mathcal O_{\mathcal A},\mathcal O_{\mathcal A}(1)\rangle^\perp
 \subset\Perf(\mathcal A)
\]
となる。比較は特殊ファイバーの標識およびArtin基底間の導来基底変換と両立する。
基底自身の特異性を含む絶対特異圏との同一視ではない。
\end{theorem}
\begin{proof}
相対対象は、$R$上一様な有限Tor振幅を持ち、各内部次数が$R$上perfectな対象に等しい。
一方向は有限$B$自由分解から従う。
逆に有限$A$自由分解を十分下で切ると、有限Tor振幅から高いsyzygyは$R$平坦となる。
$B$の有限相対対角分解をtensorして長さ五の$B$射影分解を得る。
ここで全内部次数の直和が有限階数であることは要求しない。

MFの$P$は特殊ファイバーで単射であり、各内部次数の単元minorから余核は$R$自由、$B$射影次元高々一となる。
逆に$R$平坦な有限生成$A$加群$M$は$\operatorname{pd}_BM\le5$であり、
$A$自由加群の$B$射影次元が一なので
$\operatorname{pd}_B\Omega_AM\le\max\{1,\operatorname{pd}_BM-1\}$である。
十分高いsyzygy$M'$の分解$0\to P_1\xrightarrow P P_0\to M'\to0$を取る。
$WM'=0$から$W:P_0(-3)\to P_0$は一意に$P$を通り$Q$を与える。
$PQ=W$と$P$の単射性から$QP=W$も従う。
有限graded射影は局所基底上の冪等元持上げにより自由なshiftの直和因子として扱える。
一般の相対複体も有限$A$自由複体を除けばこの状況になる。
二周期完全分解とその双対の完全性は$B$へ持ち上げて$W$の正則性を使えば得られ、
安定射、MFのhomotopy類、相対特異圏の射を同一視する。
これが完全忠実性と本質的全射性を与える。

windowに必要な内部cutoffを確かめる。
高い自由分解は上の二周期tailに一致し、二段ごとに生成次数が三増すので、
任意の$i$より小さい生成次数は十分左には現れない。
内部切断$M_{\ge i}$とその商は各次数の直和因子で有限Tor振幅を保つ。
有限内部次数の相対対象は$R(j)\otimes_R\Perf(R)$で生成される。
通常のcohomological truncationで相対部分圏が閉じることは仮定しない。
中心元の長さ一分解から
\[
 \operatorname{RHom}_B(M,B)=\operatorname{RHom}_A(M,A)(3)[-1],\qquad
 \operatorname{RHom}_A(R,A)=R(2)[-4]
\]
となる。従って$D_A=\operatorname{RHom}_A(-,A)$は相対部分圏を保ち$D_A^2=1$、
$D_A(R(j))=R(2-j)[-4]$, $D_A(A(j))=A(-j)$である。
これは相対Gorenstein parameter二の式であり、Artin環の絶対有限大域次元を用いない。

Ore chartの代数$C$と$C^e$は$R$平坦である。
特殊ファイバーの滑らかなaffine対角の有限射影分解を有限自由部分分解として持ち上げる。
各核は$R$平坦で、最後の核の剰余体還元は射影である。
冪零イデアルを越えた冪等元持上げと
$\operatorname{Tor}^{C^e}_1(C^e/\mathfrak{m}_RC^e,K)
=\operatorname{Tor}^R_1(R/\mathfrak{m}_R,K)=0$
から最後の核も射影となる。
よってchartは相対smoothで、$R$上有限Tor振幅のcoherent複体は正確にchart上perfectである。
有限affine coverのČech複体は平坦な項と一様有限振幅を持つ。
proper特殊ファイバーの有限cohomologyと冪零摂動の有限和を用いて、各twistの導来切断は$R$上perfectで基底変換と両立する。
Serre消滅と冪零filtrationにより$\mathbf R\Gamma_{\ge i}$は有限生成かつ一様有限振幅で、相対部分圏に入る。

Burke--Stevenson\cite{BS}の原定理は次数零部分の有限大域次元を仮定するため、Artin基底へ直接適用しない。
以上の内部cutoff、torsion局所化、perfect双対性、$R(j),A(j)$の像を用いて内部次数分解の直交補を比較する。
parameter二なので二つの$A$-shiftが残り、window $i=-1$は
\[
 \Perf(\mathcal A)=\langle\mathcal O_{\mathcal A},\mathcal O_{\mathcal A}(1),\mathcal B_{-1}\rangle,
 \qquad\mathcal B_{-1}\simeq D_{\rm sg,rel}(A)
\]
を与える。この順序の最後の成分をそのまま$\mathcal K$と呼ばず、
$L_{\mathcal O}L_{\mathcal O(1)}$で左へ移して
$\Perf(\mathcal A)=\langle\mathcal K,\mathcal O,\mathcal O(1)\rangle$とする。
評価複体はperfectなのでmutationも有限である。
全構成は自由分解、導来随伴、cone、dg quotient上で行われ、Morita同値へ上がる。
有限$B$分解、Ore局所化、perfect評価複体を用いるので導来基底変換と両立する。
両側で同じwindowとmutationを用いれば標識も一致する。
\end{proof}

\begin{lemma}[生成子一つの持上げによる基底変換]
$Q$は有限大域次元の可換Noether局所環、$R$はその局所Artin商とする。
次数別有限自由な$(B_Q,W_Q)$について、有限自由MFの有限直和$G_Q$があり、
$G_k$が剰余体MF圏を生成し、$A_Q=\operatorname{RHom}(G_Q,G_Q)$が$Q$上smooth properであるとする。
剰余体MF圏もproperとする。このとき
\[
 \Perf(A_Q\otimes_Q^{\mathbf L}R)\simeq\MF(B_R,W_R)^{\rm perf}.
\]
任意のArtin対象そのものの$Q$への持上げは必要ない。
\end{lemma}
\begin{proof}
有限自由graded MFの射複体は各cohomological次数が有限自由で、剰余体還元は実際の射複体である。
その剰余体複体を有限cohomologyとcontractibleな二項複体へ分け、基底を$R$へ持ち上げる。
微分の冪零摂動は有限幾何級数で消去できるので、射複体は$R$上perfectとなる。
各次数は有限自由で、摂動合成は基底の冪零次数で止まるため、一方向無限の複体でもこの操作は定義される。
$Q$上の$R$の有限射影分解により$A_Q\otimes_Q^{\mathbf L}R=A_R$を得る。
任意の$M$に対して$P=\operatorname{RHom}(G_R,M)$は$R$上perfectであり、
$A_R$のsmooth対角の有限cell分解をtensorすると$P$は$A_R$上perfectになる。
従って評価$P\otimes_{A_R}^{\mathbf L}G_R\to M$のcone $Z$は$G_R$と直交する。
導来基底変換と剰余体生成性から$Z_k=0$、perfect複体のNakayamaから$\operatorname{RHom}(Z,Z)=0$、従って$Z=0$である。
\end{proof}

\section{次数シフトを保つ楕円比較の剛性}
$E/R$は原点$o$を持つ滑らかな楕円曲線族、$L=\mathcal O_E(3o)$、
$\rho=T_{\mathcal O_E}\circ(-\otimes L)$とする。
\begin{lemma}\label{mix:rigidity}
各$R/(h^{N+1})$上で特殊ファイバーが恒等であるMorita自己同値$T$が$T\rho\simeq\rho T$を満たすなら、
$T$の同型類は恒等であり、この同型を切断間で両立して選べる。
指定された$\rho^3\simeq[2]$を保つ比較同型は、特殊ファイバーでの値を固定すれば一意である。
これは自然自己同型としてのscalarを消す主張ではない。
\end{lemma}
\begin{proof}
微小自己同値は$HH^1(E_0)=H^0(T_{E_0})\oplus H^1(\mathcal O_{E_0})$であり、
平行移動$t_a^*$とPicard tensor $A_b=-\otimes\mathcal O(b-o)$が二方向を張る。
$-\otimes L$の共役は$(a,b)\mapsto(a,b+3a)$、
球面twistの共役は$(a,b)\mapsto(a-b,b)$である。
前者は$t_a^*L=L\otimes P_{-3a}$から、後者は
$A_bT_{\mathcal O}A_b^{-1}=t_b^*T_{\mathcal O}$から従う。
よって$\rho$の接行列は
$\left(\begin{smallmatrix}-2&-1\\3&1\end{smallmatrix}\right)$で、$\det(M-I)=3$、不変部分は零である。
小拡大に沿って恒等bimoduleの持上げの差は$HH^1(E_0)\otimes I$に属し、
$\rho$と可換という条件で不変部分に入るため零となる。
$HH^0(E_0)$のscalarは同型の選択だけに残り、持上げの存在を妨げない。
前段の比較同型のscalarを任意に持ち上げて割ることで帰納を両立させる。
最後に同じfunctorの二つのrotation比較の差は$a(h)\in R^\times$で、周期性保存から$a^3=1$。
$a(0)=1$と三の可逆性から$a=1$である。
\end{proof}

\section{積とSerre traceが決める混合単元}
ここでは分離Sklyanin因子とHesse曲線の標識付き比較を入力とする。
この分離比較の証明は、現行大部稿のHesse族・Picard共役の節にあり、
一般次元の標識付き可換復元の正本はnc-cubic論文Bにある。
有限gaugeから得た比較は補題\ref{mix:finite}の$\beta$線形性により指定周期性を保つ。
分離比較をPicard tensor $J=-\otimes\mathcal O(p-o)$とし、
そのrotation比較$\eta:J\rho_N\simeq\rho_CJ$の周期性defectを
\[
 (\iota_CJ)\eta^{(3)}=\kappa(J\iota_N),\qquad
 c(0)=1,\quad c^3=\kappa^{-1}
\]
で定める。補題\ref{mix:rigidity}から有限gauge比較と$(J,c\eta)$は同型である。

内部次数二の係数$m$の実際の変分は、積と曲率を同時に微分する。
命題\ref{mix:defect}の$(R_0,R_1)$が、指定周期性を用いて読んだ
$r:M\to\rho_NM$というcharacteristic類の対象成分である。
持上げ行列を変えても境界しか変わらず、有限$A_\infty$関手のtwist恒等式を$m$で微分するとこの類の自然性が従う。
この段階で一変数二因子のHochschild類を対象上の評価で代用してはいない。
まず係数を微分しない$\mathbb Q[m]$上の比較を構成し、その後$m=\alpha uv$を代入する。

\begin{proposition}[scalar消去]\label{mix:scalar}
$f:M\to\rho_N^2M[-1]$と$r:M\to\rho_NM$について、指定周期性で
$[f^2]=a_N[r]$、$\tau_N=\operatorname{Tr}_\omega(fr)$と定める。
可換側にも$a_C,\tau_C$を同じ規約で定める。
正のrotation比較で$f\mapsto\lambda f_C,r\mapsto q r_C$なら
\[
 \lambda^2a_C=\kappa a_Nq,\qquad \lambda q\tau_C=\kappa\tau_N.
\]
従って周期性を保つ比較での混合係数$b=cq$は
\[
 b^3=\frac{q^3}{\kappa}=\frac{a_C\tau_N^2}{a_N\tau_C^2}.
\]
\end{proposition}
\begin{proof}
$f^2$の終点は四回転であり、一回転へ戻す際に周期性defectが一度現れる。
$fr$も三回転から次数一の自己射へ戻す際に同じdefectが一度現れる。
Picard tensorは同じ曲線の自己拡大とSerre traceを保つので二式を得る。
第二式から$\lambda$を消せば結論となる。
$\eta$を$u\eta$に変えても$(\lambda,q,\kappa)\mapsto(u^2\lambda,uq,u^3\kappa)$なので不変である。
必要な$h$の消去はまず$\CC((h))$で行い、最後に$R$へ戻す。
\end{proof}

\begin{lemma}[windowにおける奇射の第二行]
点MFの$Q^{\rm MF}$の第二行$\epsilon$で余核を点イデアルへ同一視する。
奇射の第一行列を$A$とし、評価列を
$0\to K\to H^0(HN)\otimes\mathcal O_E\to HN\to0$とする。
$\rho_N(f):K[1]\to N[1]$は$A$の第二行$A_2$を$K$へ制限した写像である。
評価基底が$(s_1,s_2)$なら対応切断は$A_2(-s_2,s_1)^t$である。
\end{lemma}
\begin{proof}
周期的自由分解を三つずらした二段拡張
$0\to I_o\xrightarrow{Q(3)}C(1)\oplus C\xrightarrow{P(3)}C(2)^2\to I_o(3)\to0$
の$C$成分は$\epsilon:I_o\hookrightarrow C$そのものであり、指定周期性に余分な単元を導入していない。
closed条件の第二行は$\epsilon B=-A_2P$である。
従って$A_2$は商に降りて$HN\to\mathcal O_o$を与え、
評価列と$0\to N\to\mathcal O_E\to\mathcal O_o\to0$の間の図式の中段を定める。
左端は$A_2|_K$となり、このpushout/pullbackが主張のdegree一射である。
\end{proof}

\begin{theorem}\label{mix:b}
係数を微分しない純曲率比較の混合一次係数を、Hesse規格化後に$b(h)x$と書く。
標識、次数シフト、指定三次周期性を保持する比較では
\[
 b(h)^3=1+h^6/27,\qquad b(0)=1.
\]
従って一混合族のcanonical代表は
\[
 D/3+u^3+v^3+h^2xyz/3+\alpha(1+h^6/27)^{1/3}uvx.
\]
$\alpha$の再パラメータ化を用いない。
\end{theorem}
\begin{proof}
命題\ref{mix:product}から$a_N=-h$である。
Homの座標は$\ell_N(U,V)=U_{11}(1,-1,2h)$で、境界を殺し$\ell_N(R_0,R_1)=1$となる。
可換側の$r_C=x\operatorname{id}$の座標は$U_{11}(1,-1,h)$なので、
\[
 a_C=\det A_C(1,-1,h)=-\frac{9h}{(3+h^3)^2}(1+h^6/27).
\]
混合characteristic類の評価はこの一次元方向で単射である。
実際$\operatorname{Hom}(M_C,\rho_CM_C)=H^0(L)/H^0(L(-q))$であり、$q$での評価が座標となる。
$x(q)=1$なので$x\operatorname{id}$は単元生成元である。

traceを実際のresidueで計算する。$x=1,t=z/x$のchartで
\[
 \omega=\frac{dt}{y^2+h^2t/3},\qquad\tau_0=3(x+y)-h^2z,
 \qquad\operatorname{div}\tau_0=3o
\]
とする。評価列の基底$(s_1,s_2)$と切断$u:N\to HN,v:K\to N$に対し、
二つの局所liftの差はkernel vectorの$-u/(s_1s_2)$倍となる。
同じcone符号を選ぶとtraceは$(uv/(s_1s_2))\omega$の$s_1$の零点でのresidueの和である。
この公式は評価列のČech接続写像と曲線のSerre双対のresidue表示\cite{Tate}による。

源で$N=\mathcal O(-o)$、$HN=\mathcal O(2[p])$とし、評価基底を
$s_1=(l/\tau_0)\circ t_{-p},s_2=(z/\tau_0)\circ t_{-p}$とする。
$A_2=(0,1)$なので奇射の切断は$s_1$である。
$r$の有理関数は右THCR積の規約から
\[
 r_N=\{s_1h(x-2y)/3+s_2(x+2h^2z/3)\}/z.
\]
従ってtraceは$s_1$の零因子$o+[2]p$での$(r_N/s_2)\omega$のresidueである。
$o$では$s_1/s_2$が零、$\omega=dt+O(t)dt$なのでresidueは一である。
$[2]p$では$z\ne0$なので正則で寄与は零となる。
ここで形式群法から$t([2]p)=-2h+O(h^2)$であり、分数体上で非零である。
よって$\tau_N=1$を得る。

可換側は$N_C=\mathcal O(-q)$、評価基底は$l/\tau_0,k/\tau_0$であり、
奇射の切断は$(l+rk)/\tau_0=\tau_q/((3+h^3)\tau_0)$、
$\tau_q=(3-h^3)x+(3+h^3)y+2h^2z$である。
$r_C$は$x/\tau_0$なのでtraceは
$(x/k+rx/l)\omega$の$o,p$でのresidueの和である。
曲線の陰関数微分から
\[
 l'(o)=h^2/3,\quad l'(p)=-2h^2/(3-h^3),\quad
 (\omega/dt)(o)=1,\quad(\omega/dt)(p)=3/(3-h^3).
\]
$k(o)=-h,k(p)=-2h$なので$x/k$は寄与せず、
二residueは$6/(3+h^3),-3/(3+h^3)$となる。
従って$\tau_C=3/(3+h^3)$である。
Picard tensorの同一視に用いる有理関数は自己射の共役で相殺されるため、この共通微分のtraceを保つ。
命題\ref{mix:scalar}へ代入して$b^3=1+h^6/27$を得る。

最後に一変数$m$の重み二から、全次数純曲率は$C_h(x)+mL_h(x)$である。
$h$のみに依存するHesse規格化後、固定対角群と巡回対称性、標識付きsliceの一意性から
一混合項の像は$\alpha b(h)uvx$に限られる。
分離Hesse係数は$h^2/3$であり、相対windowと有限移送により実際の残余圏変形に一致する。
定数項一の立方根は全Artin商と両立して一意である。
\end{proof}

\paragraph{六次計算への代替.}
$b=1+h^6/81+O(h^{12})$は全次数の積・trace公式の系である。
従って失われた六次gauge入力の数値をこの係数の証明として再利用する必要はない。
候補三次式への中心二変数の代数的座標変更は可能でも、比較核の収束・代数化や有限非零$h$の圏同値を証明したことにはならない。

\section{全係数を保持した二次曲率縮約}
$K=\mathbb Q[e_1,e_2,e_3,m_1,m_2,m_3]$、
$F=D/3+\sum e_ix_jx_k+\sum m_ix_i$とする。
$\pi^{ij}=\epsilon^{ijk}F_k$と置き、反復添字は$x,y,z$を走る。
\begin{lemma}\label{mix:weylsecond}
実際の周囲の完全対称化による固定化は、$h^3=0$上で
\begin{align*}
 \mu_h&=\mu+hC_1+h^2C_2,\\
 C_1(f,g)&=\tfrac12\pi^{ij}f_i g_j,\\
 C_2(f,g)&=\tfrac18\pi^{ij}\pi^{kl}f_{ik}g_{jl}
 +\tfrac1{12}\pi^{ij}(\partial_j\pi^{kl})(f_{ik}g_l-f_kg_{il})
\end{align*}
である。preferred中心の同じ固定化での像は
$F+h^2w_2$、$w_2=-2xyz/3+\sum e_i^2x_i/6$となる。
\end{lemma}
\begin{proof}
Jacobi恒等式とLeibniz則から
$[\mu,C_1]=0$、$[\mu,C_2]+[C_1,C_1]/2=0$である。
これらは多重微分作用素の全入力係数恒等式であり、現行検証器も全係数のまま確認する。
積は厳密単位的で、生成元交換子は$h\pi^{ij}$である。
$F_i$が二次以下なので所定の対称化関係式に一致する。
$\theta=\sum t_ix_i$に対して$C_1(\theta,e^{\tau\theta})=C_2(\theta,e^{\tau\theta})=0$である。
前者は反対称性、後者は$\theta$の二階微分が零で残る$\partial_j\pi^{kl}$と$t_kt_l$の縮約が零であることによる。
従って星指数は通常の指数に一致し、全ての完全対称化単項式が通常単項式へ移る。
形式PBW性からこれが実際のWeyl固定化である。
中心は
\[
 \Psi=\Omega+\sum e_i\operatorname{sym}(x_jx_k)+\sum m_ix_i
 +\frac{h^2}6\sum e_i^2x_i
 +\frac{h^4}{18}\sum(e_ix_i^2-e_je_kx_i)
\]
という自由代数証明書を持つ実際の中心\cite{Pair}を読む。
その二次までの完全対称化が表示式を与える。
\end{proof}

\begin{lemma}[再生成できる有限原始]\label{mix:w3}
\[
 v=-\frac\beta6\operatorname{Alt}(\partial_x,\partial_y,\partial_z),\qquad
 w_1=\frac1{24}\sum_{i,j}\operatorname{adj}(\operatorname{Hess}F)_{ij}\partial_i\partial_j
\]
とする。有限正規化三入力cochain $w_3$を、有理数線形系から一意のpivot規則で選んで
\begin{gather*}
 [\mu,v]=0,\qquad[F/\beta,v]=C_1,\\
 [\mu,w_3]=[v,C_1]/2,\qquad
 C_2-[F/\beta,w_3]=[\mu,w_1],\qquad
 [F/\beta,w_1]=-\det(\operatorname{Hess}F)/(8\beta)
\end{gather*}
とできる。全係数$e,m$は独立な多項式変数のままである。
\end{lemma}
\begin{proof}
最初の二式は反対称化の展開である。
$w_3$の生成手順を完全に指定する。
$[v,C_1]/2$を全入力微分指数$A$ごとに分け、通常微分係数に入力階乗の積を掛けてdivided derivativeへ直す。
三入力の正分割を辞書順の列、四入力の正分割を辞書順の行とし、
補題\ref{mix:cobar}の微分を行列にする。
左からのpivot列だけを使い、他列の係数を零に固定して有理数上の線形系を解く。
得た係数を入力階乗で割れば通常微分表示となる。

この手順を実装した現行検証器は原始のデータファイルを読まず、155入力微分型の$w_3$を毎回再生成する。
生成とは別の通常微分のLeibniz挿入で、上の二つの中央恒等式を全多項式係数について確認する。
具体的には$A\circ_iB$で外側微分指数$a_i$を内側の係数と各入力へ全て配分し、
多項係数$a_i!/(\gamma_0!\cdots\gamma_q!)$を掛ける。
括弧はarityからのGerstenhaber符号で取る。これにより任意入力に対する作用素恒等式を検証し、有限個の入力試験で代用しない。
最後の式は
$w_1(F)=\operatorname{tr}(\operatorname{adj}(\operatorname{Hess}F)\operatorname{Hess}F)/24
=\det(\operatorname{Hess}F)/8$からも直接従う。
$\beta$の冪と内部次数、入力ごとの正規化は生成過程で保たれる。
\end{proof}

\begin{theorem}\label{mix:secondcurve}
$\xi=hv+h^2(w_1+w_3)$による有限gaugeは、全$e,m$を保持して源を
\[
 \mu+\frac1\beta\left(F+h^2\left\{\frac{xyz}3-
 \frac1{12}\sum_i e_i^2x_i+\frac14e_1e_2e_3\right\}\right)
 \pmod{h^3}
\]
へ移す。このgaugeを全次数構成の初期部分に選べる。
\end{theorem}
\begin{proof}
源は$\mu+hC_1+h^2C_2+(F+h^2w_2)/\beta$である。
指数を二次まで展開し補題\ref{mix:w3}を用いると、積と正arity項は消え、
曲率は$F+h^2\{w_2+\det(\operatorname{Hess}F)/8\}$となる。
\[
 \det(\operatorname{Hess}F)=8xyz-2\sum_i e_i^2x_i+2e_1e_2e_3
\]
を代入すれば結論を得る。
指数のTaylor成分はこの切断で高々五入力である。
必要なら有限対称群で原始を平均でき、曲率を保ったまま全次数帰納を三次から再開できる。
中央係数へ線形で正規化されているので、$e_i$を$u,v$の一次式、$m_i$を二次式へ代入し$u^3+v^3$を加えても同じ式が成立する。
補題\ref{mix:finite}と定理\ref{mix:window}によりこれは実際の残余圏の標識付き可換代表である。
一般に平方混合項を含むため、係数をそのまま五変数canonical Fermat sliceの係数とは呼ばない。
\end{proof}

\section{全混合代表の六級数と中央平行移動}
preferred中心は前節の$\Psi_{e,m,h}$に固定する。
$\lambda_i^3=1,\Lambda=\prod\lambda_i$に対し
\[
 x_i\mapsto\lambda_i x_i,\quad e_i\mapsto\Lambda^{-1}\lambda_i e_i,\quad
 m_i\mapsto\lambda_i^{-1}m_i,\quad h\mapsto\Lambda h
\]
とし、置換は添字を同時置換して$h$に符号を掛ける。
Weyl対称化はこの有限群作用と可換である。
命題\ref{mix:pure}の各段の原始と曲率を有限群で平均すれば、
係数線形性、有限arity、内部次数、特殊ファイバーの標識を保った共変gaugeを得る。
有理係数性もGalois不変な平均から保たれる。

\begin{theorem}\label{mix:six}
全六非対角一次係数と九二次係数に対して、純曲率代表を
\begin{align*}
 H_h={}&D/3+h^2xyz/3+b(h)\sum_i m_ix_i+\gamma(h)\sum_i e_ix_jx_k\\
 &+U_e(h)\sum_i e_i^2x_i+V_e(h)\sum_i e_je_kx_i
   +R_e(h)\sum_i e_i^3+S_e(h)e_1e_2e_3+T_e(h)\sum_i e_im_i
\end{align*}
と選べる。ここで$b^3=1+h^6/27,b(0)=1$であり、
\begin{gather*}
 \gamma\in1+h^6\mathbb Q[[h^6]],\quad
 U_e,S_e\in h^2\mathbb Q[[h^6]],\quad
 V_e,T_e\in h^4\mathbb Q[[h^6]],\quad R_e\in h^6\mathbb Q[[h^6]],\\
 U_e=-h^2/12+O(h^8),\qquad S_e=h^2/4+O(h^8).
\end{gather*}
六つの未知級数の全値や、この代表の無標識での一意性は主張しない。
\end{theorem}
\begin{proof}
$\Lambda=1$の対角群と巡回群の不変な重み三単項式を列挙すると、
\begin{gather*}
 D,\ xyz,\ \sum e_ix_jx_k,\ \sum e_ix_i^2,\ \sum e_i^2x_i,\\
 \sum e_je_kx_i,\ \sum e_i^3,\ e_1e_2e_3,\ \sum m_ix_i,\ \sum e_im_i
\end{gather*}
の十型となる。
全て置換対称なので係数は$h$の偶級数である。
$D/3$係数を$A$、$\sum e_ix_i^2$係数を$E$とすると
$x_i=A^{-1/3}X_i-(E/A)e_i$で前者を一、後者を零にできる。
$e=0$で三つの独立$m_i$の指標方向はそれぞれ$x_i$のみで、巡回対称性から共通係数となる。
一混合項で定理\ref{mix:b}を用いるとその係数は$b$、純$xyz$は$h^2/3$である。
重みからこれらの係数がさらに$e,m$に依存する余地はない。
残る対角群の$h\mapsto\Lambda h$と置換偶性を合わせると表示の合同条件を得る。
二次初期gaugeを定理\ref{mix:secondcurve}のものに固定すれば、最終線形規格化は
共通拡大$1+O(h^6)$と平行移動$O(h^4)$なので二つの二次係数を変えない。
\end{proof}

\begin{proposition}[全二十一係数と六自由方向]\label{mix:translation}
$x_i$の対角線形項を$L_i$、非対角項を$e_i$、二次項を$m_i$とする。
$X_i=x_i+L_i/2$とし、
\[
 n_i=m_i-L_i^2/4-(e_jL_k+e_kL_j)/2,\qquad
 C_{\rm cl}=\sum_i L_i^3/12+\sum_i e_iL_jL_k/4-\sum_i m_iL_i/2
\]
と置く。共変規格化中心を用いた全混合の代表は
$H_h(X;e,n)+u^3+v^3+C_{\rm cl}$である。
$e=0$の十五係数部分族では追加六級数を用いず
$D(X)/3+h^2X_1X_2X_3/3+b\sum n_iX_i+u^3+v^3+C_{\rm cl}$となる。
全二十二変数源は、$h$を固定する接写像恒等の形式基底変換により、
$L=0$の十六変数源と六次元自由円板との積に分解する。
\end{proposition}
\begin{proof}
中央平行移動は交換子を変えず、実際の中心の共変性により表示の$n,C_{\rm cl}$を与える。
全gaugeは係数を微分せず正規化されているので、追加中央曲率をそのまま残す。
$g=u^3+v^3+C_{\rm cl}$は一意な$M\equiv I_2$により$g\circ M=U^3+V^3$へ移る。
実際$\mathfrak{gl}_2$の軌道微分は
$(\alpha,\beta,\gamma,\delta)\mapsto3\alpha u^3+3\beta u^2v+3\gamma uv^2+3\delta v^3$で可逆である。
混合係数イデアルを$J$とすると$C_{\rm cl}\in J^2$なので$M-I\in J^2$。
従って$(L,e,m,h)\mapsto(L,e\circ M,n\circ M,h)$の接写像は恒等で、形式自己同型となる。
新座標の中心対に$L$は現れず、それを自由座標として残せる。
二座標置換は次数付き、特殊ファイバーで恒等、中心の式そのものを保つので指定周期性も保つ。
PBW定数零など別の中心規格化を用いる場合には、その中央scalar差を同じ向きで二元三次式へ移す必要がある。
\end{proof}

\section{標識付き平面incidenceへの同定}
標的のsquarefree Fermat切片を
\[
 F_{s,r}=F_0+sxyz+\sum_i(P_iu+Q_iv)x_jx_k+\sum_iT_ix_iuv,
 \qquad r=(P,Q,T)
\]
とする。$u=\sum\lambda_ix_i,v=\sum\mu_ix_i$で制限し、三変数marked Hesse規格化した係数を$\sigma$と書く。
\begin{theorem}\label{mix:incidence}
$L=0$の実際の全形式対応は、$h$を固定する$A=\CC[[s,r]]$代数同型により
\[
 B_{16}\simeq A[[\lambda,\mu,h]]/(\sigma(s,r,\lambda,\mu)-h^2/3)
\]
と表示される。全族はこの源に六つの自由変数を加えたものとなる。
原点で$\lambda_i=a_i/3+O(2),\mu_i=c_i/3+O(2)$、$r=(p,q,d)+O(2)$である。
従って平面incidence稿の臨界幾何を実際の形式対応へ適用できる。
\end{theorem}
\begin{proof}
定理\ref{mix:six}の可換代表を$G\equiv I_5$でcanonical Fermat切片へ移す。
元の平面$W_0=(u=v=0)$上の制限は正確に$D/3+h^2xyz/3$である。
新座標の平面は$G^{-1}W_0$であり、切片の制限式はこのHesse式に近恒等三変数線形変換を施したものとなる。
marked Hesse規格化の一意性から$\sigma=h^2/3$を得る。
一次shearは$u\mapsto u-\sum a_ix_i/3,v\mapsto v-\sum c_ix_i/3$なので、表示の接写像となる。
$\partial_s\sigma(0)=1$より右辺は正則十六次元環であり、
$(r,\lambda,\mu,h)$の接写像は可逆である。
形式逆関数定理で同型を得る。全族は命題\ref{mix:translation}から従う。
この同型は形式座標のもので、比較核の有限$h$での代数化を述べていない。
\end{proof}

\begin{corollary}[接階数と正確な臨界イデアル]
分離軸上のcanonical接写像は
$\partial_{p_i}\mapsto\gamma[ux_jx_k]$、
$\partial_{q_i}\mapsto\gamma[vx_jx_k]$、$\partial_{d_i}\mapsto b[x_iuv]$であり、
$L,a,c$の十二方向は零である。従って混合接階数は九、
量子方向も含めた十次minorは$(2h/3)b^3\gamma^6$である。
原点では階数九、分離軸の形式分数体上では十となる。
$L=0$の実際の源における相対臨界イデアルは正確に
$(h,p_1,p_2,p_3,q_1,q_2,q_3)$である。
全族では六つの自由な$L$方向が加わる。
\end{corollary}
\begin{proof}
分離Jacobian環の次数三基底は$xyz,x_iuv,ux_jx_k,vx_jx_k$である。
定理\ref{mix:six}を係数の一次まで取れば接写像とminorを得る。
定理\ref{mix:incidence}で相対臨界イデアルは$(h,\partial_\lambda\sigma,\partial_\mu\sigma)$である。
$h=0$では源の可換代表は古典三次式そのもので、$W_0$上の制限はFermat式、
二つの正規二次類は$\sum p_ix_jx_k,\sum q_ix_jx_k$となる。
内在的臨界条件の証明の完全対合は、これら六係数に対して可逆な線形写像なので、
零スキームのイデアルは$(p,q)$そのものである。
形式座標変更のJacobianは可逆だから、このイデアル同定は集合論だけでなくスキーム論的に成立する。
\end{proof}

\begin{thebibliography}{9}
\bibitem{Takahashi} A.~Takahashi, \emph{Matrix Factorizations and Representations of Quivers I},
\url{https://arxiv.org/abs/math/0506347}, Theorem 3.1 and Remark 3.6.
\bibitem{BFK} M.~Ballard, D.~Favero and L.~Katzarkov,
\emph{A category of kernels for equivariant factorizations and its implications for Hodge theory},
Publ. Math.\ IHES 120 (2014), 1--111, Theorem 5.39, Definition 6.24 and Lemma 6.25,
\url{https://doi.org/10.1007/s10240-013-0059-9}.
\bibitem{Pair} 中心三次対の現行論文、
\href{https://github.com/uedakazushi/nc-cubic2/blob/main/papers/nc_cubic_threefold_moduli.tex}
{非可換三次三様体の17次元モジュライ族}、全混合PBW性・中心式。
\bibitem{BS} J.~Burke and G.~Stevenson, \emph{The derived category of a graded Gorenstein ring},
\url{https://arxiv.org/abs/1507.00830}, Lemma 3.10, Propositions 4.18 and 5.8, Theorem 6.4.
\bibitem{Tate} J.~Tate, \emph{Residues of differentials on curves}, Ann. Sci.\ ENS (4) 1 (1968), 149--159,
\url{https://doi.org/10.24033/asens.1162}.
\bibitem{BW} S.~Barmeier and Z.~Wang, \emph{Deformations of path algebras of quivers with relations}, \url{https://arxiv.org/abs/2002.10001},
Theorem 10.3, Lemma 10.9 (v4).
\end{thebibliography}
\end{document}
